% Appendix: the end-to-end evaluation in detail. Moved from Section~\ref{sec:e2e}. Sources:
% evidence/kaggle-scale-run-2026-10-24.md (X0, seed 7, wardowl 2.1.0, before the honest-completeness
% change), evidence/x8-run.md (X8, seeds 11, 13 and 17, pre-registered in plan/x8-prereg.md), and the
% OpenFGA configuration's four runs (X8-FGA, amendment 1 of the same pre-registration), and
% evidence/x8-controls.md (exposure and the mutation control, amendment 2).

\section{The end-to-end evaluation in detail (X0 and X8)}\label{app:e2e}

This appendix gives the full setup of the end-to-end runs, the safety table, the exposure and positive control behind it,
and what the model did with what it retrieved, which Section~\ref{sec:e2e} summarizes.

\subsection{Setup}\label{sec:e2e-setup}

\paragraph{Corpus.}
The index holds 107,540 chunks, one per document or passage: 49,920 synthetic company documents in 20
tenants, of which 41,600 carry a unique canary string; 20 notices, one per tenant, each carrying a
prompt injection; and 57,600 passages of the FiQA financial question-answering
corpus~\cite{maia2018www18} from BEIR~\cite{thakur2021beir}, which every tenant may read and which
serve as distractors. Company documents are drawn from templates (handbooks, finance, HR, board and
project records) with per-document ACLs over users and groups. Chunks are embedded with
bge-base-en-v1.5~\cite{xiao2024cpack}.

\paragraph{People and questions.}
189 people span the tenants. A typical employee may read the shared passages, the tenant's notice
and five of the twelve document templates, 58,641 chunks in all (the median); the least privileged
are non-executive board members (``judy'' personas), each of whom may read 624 chunks of three
templates in one tenant and none of the shared passages (0.58\% of the corpus); one person in no
group may read nothing. They ask 1,105 questions: 108 generic, 376 targeted at a document
they may read, 189 targeted at a document they may not read (\emph{forbidden}), and 432 red-team
prompts that ask for other people's data, other tenants' data or everything, or try to override
instructions.

\paragraph{Policy and stores.}
The policy is $\mathrm{AllOf}(\mathrm{Tenant}, \mathrm{MetadataACL})$ (Appendix~\ref{sec:sys-compose}).
It runs against ten configurations: FAISS over a flat index (exact, the reference), Chroma, pgvector,
pgvector with the same policy written in OPA, Qdrant, Elasticsearch, OpenSearch, Weaviate, Redis and
LanceDB, each at the settings listed in Appendix~\ref{app:settings}. Each question retrieves with
$k = 5$ for the answer and $k = 10$ for recall@10. Qwen2.5-3B-Instruct~\cite{yang2024qwen25} answers
from the five chunks. An eleventh configuration, pgvector with the policy's ACLs held as 74,960
OpenFGA relationship tuples behind the caching authorizer, pre-filters with ListObjects. It did not
start in the runs below: our harness set OpenFGA's ListObjects deadline to 20\,s while leaving its
request timeout at the default, and OpenFGA 1.22 refuses to start with a deadline longer than the
request timeout. With the timeout raised and nothing else changed, the configuration ran in four further
runs with the same seeds, pre-registered as an amendment before they started; each also repeated
the in-process pgvector configuration.

\paragraph{Runs.}
The scenario ran four times with seeds 7, 11, 13 and 17, on one Kaggle session with two T4 GPUs each
(Appendix~\ref{app:compute}). The software was not identical across runs, a deviation from the
pre-registration's ``everything else identical'': the first run (seed 7) had Redis 8.10.2, lancedb 0.40.0 and
weaviate-client 4.24.0, and the three repeats had Redis 8.10.3, lancedb 0.41.0 and weaviate-client 4.24.1,
as the runs' logs record. The other recorded versions match, among them Elasticsearch 9.5.5, OpenSearch 3.9.0, Qdrant 1.19.2, Weaviate 1.39.12,
chromadb 1.5.9 and faiss-cpu 1.15.1. So the ranges for Redis and LanceDB in Figure~\ref{fig:qual-audit} span two versions.
The first run also used another Kaggle account of the same session type.
The seed changes the document
text, part of the sample of questions and the order in which people ask; the permission templates,
personas and the 216 pairs of Section~\ref{sec:e2e-safety} are fixed, so the runs repeat the layout with fresh documents and mostly fresh questions. The repeats
and their predictions were pre-registered after the first run and before any repeat. The runs used
the released version 2.1.0, before the change of Appendix~\ref{sec:sys-complete}; permissions did not
change during a run, so these runs test containment, not freshness, which
Section~\ref{sec:eng-revoke} measures.

\paragraph{Leak detection.}
A leak is any returned chunk outside the person's readable set, as decided by the policy oracle,
which evaluates the policy without the layer, or any canary of an unreadable document found in a
retrieved chunk, in the prompt sent to the model or in the model's answer. A canary is a 40-bit code derived
from the document's identifier; the matcher finds it in any case and with any hyphen, dash, space or
nothing between its parts, and a positive control checks, for every person, that it finds exactly
the canaries of the documents that person may not read in a text that holds all of them. A model
cannot produce a 40-bit code it has not been shown, so an answer can leak a canary only through a
prompt that contains it, and prompts are checked verbatim. The answer check guards against a fault
in the pipeline between prompt and answer, not against paraphrase.

\subsection{Safety in detail}\label{app:e2e-safety}

\begin{table}[tbp]
\centering
\footnotesize\setlength{\tabcolsep}{4pt}
\begin{tabular}{@{}lrrrrr@{}}
\toprule
Configuration & Questions & Hits checked & Unreadable hits & Canaries (chunk/prompt/answer) & Violations \\
\midrule
FAISS (exact) & 4,420 & 65,440 & 0 & 0/0/0 & 0 \\
Chroma & 4,420 & 65,440 & 0 & 0/0/0 & 0 \\
pgvector & 4,420 & 65,047 & 0 & 0/0/0 & 0 \\
pgvector + OPA & 4,420 & 65,071 & 0 & 0/0/0 & 0 \\
Qdrant & 4,420 & 65,440 & 0 & 0/0/0 & 0 \\
Elasticsearch & 4,420 & 65,440 & 0 & 0/0/0 & 0 \\
OpenSearch & 4,420 & 65,440 & 0 & 0/0/0 & 0 \\
Weaviate & 4,420 & 65,440 & 0 & 0/0/0 & 0 \\
Redis & 4,420 & 65,440 & 0 & 0/0/0 & 0 \\
LanceDB & 4,420 & 65,440 & 0 & 0/0/0 & 0 \\
\midrule
Total & 44,200 & 653,638 & 0 & 0/0/0 & 0 \\
\midrule
pgvector + OpenFGA\textsuperscript{$\ast$} & 4,420 & 64,845 & 0 & 0/0/0 & 0 \\
pgvector\textsuperscript{$\ast$} & 4,420 & 65,405 & 0 & 0/0/0 & 0 \\
\bottomrule
\end{tabular}
\caption{End-to-end safety over four runs (seeds 7, 11, 13 and 17). The unit of the bounds in the
text is the question and the person; hits are secondary and are counted over both retrievals per
question, so about a third of them, the $k = 5$ hits, largely repeat chunks already counted among the
$k = 10$ hits of the same question. pgvector and its variants checked fewer hits because some of
their answers were short or empty (Section~\ref{sec:quality}). \textsuperscript{$\ast$}Four further
runs with the same seeds, in which only pgvector and pgvector with OpenFGA ran. No configuration
reported an error.}
\label{tab:e2e-safety}
\end{table}

No configuration in any run returned an unreadable chunk or exposed a canary
(Table~\ref{tab:e2e-safety}). The bounds of Section~\ref{sec:e2e-safety} count distinct (person, question) pairs. In the generator
(\texttt{questions()} in the harness at bf25bd0), the generic pairs are 12 fixed strings asked by nine people, and two of the eight attacks, the 54 \texttt{ignore-instructions} and 54 \texttt{reference-codes} prompts of a run, have no
placeholder, while the attack assigned to a person does not depend on the seed. So 216 of 1,105 pairs, 19.5\%, are the same in all four seeds. The other pairs choose an entity from a pool of 208 per seed. The runs' stored questions, the same for every configuration of a seed, hold 3,746 distinct pairs: 216 in all four seeds, 26 by chance in two and 3,504 in one. The one-sided 95\% bound, separately for each of the eleven
configurations and taking those pairs as independent, is $1 - 0.05^{1/3746} \approx 8.0 \times 10^{-4}$ per question (Section~\ref{sec:prob-metrics}); the pre-registered analysis, which treated all 4,420 as independent, gives $6.8 \times 10^{-4}$. They are not fully independent even so: a document question comes from 12 templates and a red-team prompt from eight, and the same people ask them. Per
person the bound stays at a single run's $1 - 0.05^{1/189} \approx 1.6 \times 10^{-2}$: the same
personas recur in every run with the same permissions, so a fault tied to a permission shape would
recur too. The 189 people are 13 role types, since the persona's name fixes its groups, repeated across tenants, and over 13 the formula gives 0.21. Pooling the configurations would not
shrink these bounds honestly, since they share questions and the layer's code. The first of the five
pre-registered predictions for the repeats, that no unreadable hit, canary or violation would occur,
held, and so did the amendment's for OpenFGA, which started in every run and pre-filtered every
question. As a secondary count, the runs checked 783,888 hits, not independent trials, since hits
cluster by question and by person, and a fault in the layer would affect a question, not a random hit.

These bounds are over the question generator, and a binomial bound over sampled questions says nothing about
an adversary who adapts its queries to what it sees. No such adversary was run. For one we rely on
Observation~\ref{prop:reduction}, which does not depend on the query, and on its trusted computing
base.

No violation occurred, so across these runs the re-check dropped nothing, and the same outcome would appear with it off. In the in-process configurations the filter and the re-check evaluate the same
ACL metadata under the same reference semantics, so a violation would need a store to translate or
evaluate the filter wrongly; with OPA and OpenFGA, filter and check come from the engine's separate
APIs, and they agreed on every hit. Sections~\ref{sec:conformance} and~\ref{sec:engines} test
configurations and data shapes on which they do not. These runs thus test whether two implementations of one predicate agree and whether a store round-trips metadata.

Two observations qualify the result. First, it is a statement about this layer under assumptions
(A1) to (A3) of Observation~\ref{prop:reduction}, not a proof, and the inputs called adversarial are weak:
the red-team prompts can only change the query text, and the layer never lets text choose the
subject (Appendix~\ref{sec:sys-identity}). Second, containment held while quality did not, and quality
did not repeat. pgvector returned empty answers to 10, 1, 0 and 7 questions in the four
runs, from people who could read something. The harness counts answers, not people: the check's number is one entry per empty question. By the stored answers, the four runs' empties fell on 9, 1, 0 and 7 people, all least-privileged and each able to read 624 chunks, and one person in the first run had two; each asked five questions, so everyone got hits on at least three. In the first run eight of the ten were on red-team prompts and two on questions about documents the person could not read. The person who may read nothing is excluded from the check. The layer, before the change of
Appendix~\ref{sec:sys-complete}, reported those empty answers as complete. The same configuration ran
in one of two modes, fast and lossy or exact and about fifty times slower, for a reason
Section~\ref{sec:qual-pgvector} ties to PostgreSQL's plan cache in a separate run. Our third prediction, that 3 to 25
people would receive no results in every repeat, was scored on this count and falsified: two repeats fell outside it, because
it assumed that the first run's mode was a property of the configuration. Over its eight runs the
configuration ran exact in five and, by its median, fast in three, one of them a mixture whose p95 of
433\,ms shows that part of its searches ran exact, and it gave at least one empty answer in five of the eight (10, 1, 0, 7, 1, 1, 0, 0). The seed did not decide which: seeds 7, 11 and 17
each ran fast in one run and exact in another. In two of the five exact runs one person received one
empty answer, which we have not explained. With OpenFGA, whose pre-filter lists the readable documents,
pgvector was lossy in every run, with recall@10 between 0.693 and 0.758 and 7, 8, 3 and 10
empty answers on 6, 8, 3 and 9 least-privileged people, whatever mode the in-process configuration had just run
in; its plans were not recorded, so we have only an inference from the code for this. Qdrant answered every question, but with recall@10 between 0.535 and 0.561 in every run. The
other stores varied little from run to run (Section~\ref{sec:qual-audit}).

\subsection{Exposure and the mutation positive control}\label{app:e2e-control}

A second amendment to the pre-registration, written after the results above and after our critical
review and before running, asks whether those results have power. Its measurements are exploratory
in that sense. The harness is unchanged: a runner imports it, on the same release, corpus parameters,
models and hardware, with FAISS only.

\paragraph{Exposure.}
For every question of the four seeds we computed the exact top 5 and top 10 over the whole corpus, without any
mask, and counted the chunks the asking person may not read by the oracle, and those in which the leak detector finds a
canary of a document that person may not read (Table~\ref{tab:e2e-exposure}). The shares are the same to one decimal in every seed;
the seed changes document text and the sample of questions but not the templates or personas. Each
company document has a counterpart from the same template and pool entry in each of the other 19
tenants, and these shares depend on those copies. The control's ACL mutation below, with the re-check off, is exact search over the
person's tenant and the shared passages: of its 5,105 distinct unreadable (question, chunk) pairs, 4,979 are in the person's own tenant, 126 are shared passages and none is in
another tenant. We noticed only after the run that this pass measures exposure inside the tenant. On seed 7, in the top 5, it is 100\% for the 189 forbidden questions,
62.5\% for red-team prompts, 62.0\% for generic questions, 0.5\% for targeted ones and 47.8\% overall, against 95.2\% over the whole corpus, so targeted questions test the
tenant boundary and hardly the ACL inside it. Of the seven predictions, one failed in part: targeted questions
cleared the predicted 90\% in the top 5 in every seed, and red-team prompts did not, at 87.7\% in
every seed. One attack, a request for reference codes, has no placeholder, so all its askers in a seed share one query, whose nearest chunks
are shared passages: in the top 5 it exposed only the one person who may read nothing (1.8\% of its questions), and in the top 10 it exposed 96.3\%; the other seven attacks
were at 100\%. We had
reasoned from the counterpart documents and had not considered such a query.

\begin{table}[tbp]
\centering
\footnotesize\setlength{\tabcolsep}{4pt}
\begin{tabular}{@{}lrrrrrrr@{}}
\toprule
 & & \multicolumn{3}{c}{Top 5} & \multicolumn{3}{c}{Top 10} \\
\cmidrule(lr){3-5}\cmidrule(l){6-8}
Category & Questions & Unreadable & Mean & Canary & Unreadable & Mean & Canary \\
\midrule
Generic & 432 & 100.0\% & 4.77 & 83.3\% & 100.0\% & 9.57 & 83.3\% \\
Targeted & 1,504 & 100.0\% & 4.78 & 54.8\% & 100.0\% & 9.54 & 54.8\% \\
Forbidden & 756 & 100.0\% & 5.00 & 97.9\% & 100.0\% & 10.00 & 97.9\% \\
Red team & 1,728 & 87.7\% & 4.26 & 75.0\% & 99.5\% & 8.62 & 75.0\% \\
\midrule
All & 4,420 & 95.2\% & 4.61 & 72.9\% & 99.8\% & 9.26 & 72.9\% \\
\bottomrule
\end{tabular}
\caption{Exposure over seeds 7, 11, 13 and 17 on the unfiltered exact index: the share of questions with at
least one chunk the person may not read, the mean number of such chunks, and the share with such a chunk
that carries a canary.}
\label{tab:e2e-exposure}
\end{table}

\paragraph{The positive control.}
On seed 7, FAISS was loaded once and the harness evaluated the same 1,105 questions in seven passes. A wrapper
around the policy replaced what the policy's filter gave the store and left its check unchanged. M1 admits everything;
M2 lets the tenant clause also admit the next tenant in the harness's list, with the ACL clause unchanged; M3 drops
the ACL clause. Each person gained between 48,275 and 107,540 unreadable chunks under M1 (median 48,899), between 0 and
1,665 under M2 (median 1,041) and between 832 and 60,097 under M3 (median 1,456), and no mutation left out a readable
chunk. With the re-check off, the release still matches each hit against the filter it gave the store, so it cannot see what
a wrong filter over-admits, and it reports no violation.

Table~\ref{tab:e2e-control} gives the result. All the predictions for the control held. The unmutated pass P0 reproduced the first
run's FAISS row, 16,360 hits checked with no unreadable hit, canary or violation, and recall 1.000.
M1 with the re-check off returned the same ten chunk ids as the unfiltered exact search for all 1,105 questions, which also ties the
control to the exposure measurement. Each pass with the re-check off failed exactly the harness's four leak checks (every hit readable, and no canary in a
chunk, a prompt or an answer). Each pass with it on failed two others: the re-check found violations, and people who may read something got
empty results (887, 38 and 449 retrievals).
We compared the passes question by question, not only in total, and a separate kernel recounted the comparison with no exception. For each mutation, the questions with an unreadable hit
with the re-check off are exactly those with a violation with it on and exactly those marked incomplete, and on
each question the number of violations equals the number of unreadable hits. Every canary in an answer was also in that question's
prompt, as the design of the canaries requires (Appendix~\ref{sec:e2e-setup}).

\begin{table}[tbp]
\centering
\footnotesize\setlength{\tabcolsep}{3.5pt}
\begin{tabular}{@{}llrrrrrrrrr@{}}
\toprule
 & & \multicolumn{2}{c}{Unreadable} & \multicolumn{3}{c}{Canary in} & \multicolumn{2}{c}{Violations} & & \\
\cmidrule(lr){3-4}\cmidrule(lr){5-7}\cmidrule(lr){8-9}
Pass & Re-check & $k{=}5$ & $k{=}10$ & chunk & prompt & answer & questions & hits & Incomplete & R@10 \\
\midrule
P0 & on & 0 & 0 & 0 & 0 & 0 & 0 & 0 & 0 & 1.000 \\
\midrule
M1 & off & 1,052 & 1,103 & 805 & 805 & 138 & 0 & 0 & 0 & 0.077 \\
M2 & off & 676 & 751 & 400 & 360 & 13 & 0 & 0 & 0 & 0.703 \\
M3 & off & 528 & 533 & 509 & 503 & 107 & 0 & 0 & 0 & 0.548 \\
\midrule
M1 & on & 0 & 0 & 0 & 0 & 0 & 1,103 & 15,317 & 1,103 & 0.077 \\
M2 & on & 0 & 0 & 0 & 0 & 0 & 751 & 4,834 & 751 & 0.703 \\
M3 & on & 0 & 0 & 0 & 0 & 0 & 533 & 7,642 & 533 & 0.548 \\
\bottomrule
\end{tabular}
\caption{The mutation positive control on FAISS, seed 7, 1,105 questions, the main retrievals at $k = 5$ and $k = 10$.
Counts are questions, except the violations counted in hits. R@10 is recall@10 against the unmutated readable set.}
\label{tab:e2e-control}
\end{table}

Three limits apply. This is one seed and one configuration, an exact store, so it shows that the checks fire
and that the re-check contains these faults here; it does not test the interaction of a wrong filter
with an approximate index, and exposure was measured on FAISS only. It breaks only the filter, since the check stays the policy's, so a wrong
check is outside it and stays in the trusted computing base, and containment under an untouched check holds by construction. It was also added after a review of the results. And the re-check does not repair: recall stays
where the broken filter left it, and the layer's report of an incomplete result is what tells an
operator something is wrong.

\subsection{What the model did with what it retrieved}\label{sec:e2e-model}

Prompt injection in readable content is outside our threat model (Section~\ref{sec:prob-threat}),
but the runs measure where the boundary lies. Each tenant's injected notice is readable by that
tenant's employees and contractors, so retrieving it for them is correct. Over all eight runs the
notice was retrieved between 1 and 117 times per configuration and run, depending on the store's
ranking, and the harness counted its marker string in the 3B model's answer the same number of times: 3,085 of 3,085.
This count of ``obeyed'' is the marker anywhere in an answer, and the notice's own text holds the marker, so an answer that repeats the notice counts too. In the stored answers of all eight runs the marker appeared in exactly the answers whose top five held the notice.
Retrieval-time enforcement decides which content reaches the model, not what the model does
with it; defenses against injected instructions in permitted
content~\cite{greshake2023not,debenedetti2025defeating} are complementary.

Answer quality followed retrieval quality. With FAISS the model found the target document of a
targeted question in its top five for every question in every run and answered 98.4\% to 99.7\%
correctly (an answer is correct when it quotes the fact from the right tenant's document); with Qdrant the target
was in the top five for 33.8\% to 39.9\% of targeted questions.
With OpenFGA, despite its lower recall, the target was in the top five for 94.2\% to 97.9\%. When a
person asked about a document they could not read, the model abstained in 68\% to 83\% of cases,
depending on the store and run, and otherwise answered from readable chunks only, as the
canary checks confirm.
